\documentclass[../main.tex]{subfiles}
\begin{document}
%==========================================TIÊU ĐỀ TẬP========================================
\begin{table}[h]
  \begin{adjustwidth}{-1cm}{}
    \begin{tabular}{>{\centering\arraybackslash}p{6cm}|>{\raggedright\arraybackslash}p{12.5cm}}
      \multirow{5}{6cm}
      % Cột bên trái: ÉPISODE và số tập
      {\\[-40pt]
        \flaregothic\fontsize{20pt}{20pt}\selectfont \centering
        \color{eptype}ÉPISODE\color{black}\\
        \burbank\fontsize{72pt}{72pt}\selectfont
        \color{epnum}108\color{black}\\[6pt]
        \cafeteria\fontsize{20pt}{20pt}\selectfont\color{epnum}(8-11)\color{black}
        \\[18pt]
        \flaregothic\fontsize{18pt}{18pt}\selectfont
        \color{epattr}FIN DE TOME !\\
        \color{black}
      }
       &
      % Dòng thứ nhất: tên tập tiếng Nhật
      {\fotskip\fontsize{18pt}{18pt}\selectfont \ruby{形}{かたち}から\ruby{入}{はい}る\ruby{タ}{た}\ruby{イ}{い}\ruby{プ}{ぷ}の\ruby{名}{めい}\ruby{探}{たん}\ruby{偵}{てい}
      } \\[6pt]
      % Dòng thứ hai: tên tập tiếng Anh
       & {\cafeteria\fontsize{18pt}{18pt}\selectfont Looking the Part}                                                                            \\[4pt]
      % Dòng thứ ba: tên tập tiếng Việt
       & {\vnmsans\fontsize{18pt}{18pt}\selectfont Thám tử "lừng danh" Shion}                                                                     \\[8pt]
      % Dòng thứ tư: Nhân vật xuất hiện
       & {\caviar{\fontsize{14pt}{14pt}\selectfont Nhân vật: \emp{kuon1} \emp{ryuuhou1} \emp{achan} \emp{choco} \emp{shion} \emp{subaru} \emp{pekora} \emp{towa} }}
      \\[8pt]
      % Dòng cuối cùng: Thời gian diễn ra của tập (thường sẽ là ngày phát hành)
       & {\continuum\fontsize{16pt}{16pt}\selectfont Ngày phát hành: 30/05/2021}                                                                  \\[18pt]
    \end{tabular}
  \end{adjustwidth}
\end{table}
%=========================================TRUYỆN CHÍNH========================================
\color{black}

\txtbx[2]
{{\color{vol} \uline{Tóm tắt:}} Shion đóng vai ``thám tử lừng danh'' giải vụ mất cà rốt của Pekora, đuổi khách Choco và sợ máu khi thấy A-chan ``gục ngã''. Trợ lý thám tử của cô là Josephine - một con chó biết đi bằng hai chân và có tư duy logic, nhưng chính nó lại tố cáo chủ nhân là thủ phạm. Subaru - cảnh sát dạo - lập tức còng số 8 và giải Shion về đồn. Hóa ra toàn bộ vụ án là trò đùa của Towa, người đã dựng lên kịch bản từ đầu để ``tống khứ'' thám tử vô dụng. Họ cùng nhau tới đồn cảnh sát, thương lượng với cô nàng Vịt và bảo lãnh được nàng Phù thủy về.}

\pov{Hikawa Kuon}

Ngày cuối cùng của tháng Năm đang cận kề. Không khí mùa hè đã len lỏi qua từng con phố của thành phố Điện tử, mang theo hơi nóng oi nồng xen lẫn những cơn gió ấm phả qua các tòa nhà cao tầng. Cái nắng không còn dịu nhẹ như đầu xuân nữa, đã trở nên gay gắt hơn, báo hiệu mùa hè thực sự đã đến rất gần. Dù vậy, cái nóng vẫn chưa đủ để ép chúng tôi phải bật điều hòa - ít nhất là vào lúc này.

Thế nhưng, giữa bầu không khí rực rỡ ấy, văn phòng \hll\ lại chìm trong bóng tối một cách kỳ lạ. Không phải vì mất điện hay ai đó quên thanh toán hóa đơn, mà nơi này đang chuẩn bị cho một màn ra mắt vô cùng hoành tráng. Thực ra, là màn ra mắt của một người rất quen thuộc, nhưng ở một ``cương vị'' hoàn toàn mới.

Một ánh đèn duy nhất bật lên, rọi xuống ngay giữa phòng, tạo nên khung cảnh đầy kịch tính. Và trong vòng tròn ánh sáng ấy, một nhân vật bí ẩn từ từ lộ diện: ngồi khoanh tay đầy tự tin trên chiếc ghế xoay, môi cắn một chiếc tẩu thuốc - dĩ nhiên không phải thuốc thật, mà là loại thổi ra bong bóng.

Đó chính là ``Thám tử lừng danh'' Murasaki Shion.

\img{8-11a}

Gọi là ``thám tử'' cho oai vậy thôi, chứ thực chất em ấy chỉ khoác lên mình một phong cách có phần nửa vời: chiếc kính một tròng mà tôi mượn từ bố, mũ phớt kiểu quý ông ``xin'' được từ ông ngoại, cộng với bộ trang phục phù thủy quen thuộc mà ai cũng đã quá rõ. Vừa không liên quan chút nào, vừa tạo ra một phong thái\dots\ rất Murasaki Shion.

Ngồi ở góc phòng, Ryuuhou đang che miệng, cố nén tiếng cười lớn để không làm hỏng không khí hào hùng của nàng thám tử mới vào nghề. Cô chạm nhẹ vào cánh tay tôi, thì thầm với giọng trêu chọc: ``Anh nhìn Shion-san kìa, như đang diễn một vở kịch độc diễn ở câu lạc bộ vậy. Bộ đồ phù thủy đi với kính một tròng\dots\ hợp đến mức không thể chê được luôn.''

``Thêm một bị cáo nữa đã lộ diện!'' Shion hất cằm đầy kiêu hãnh, chỉ thẳng ngón tay về phía\dots\ không khí. Giọng nói tràn đầy phấn khích, còn ánh mắt lại toát lên vẻ nghiêm túc đến khó tin.

Tôi đứng một bên, khoanh tay quan sát cảnh tượng, trong khi Fukuyasu - người tôi đặc biệt mời đến hôm nay - và Towa đang trợn mắt nhìn nhau, chưa hiểu chuyện quái gì đang diễn ra. Ryuuhou thì rón rén di chuyển ra gần hơn, như thể đang xem một bộ phim trinh thám và không muốn bỏ lỡ cảnh cao trào nào. Cô thậm chí còn rút điện thoại ra, lén bật chế độ quay màn hình.

\img{8-11b}

Cậu bạn tên An hất cằm về phía tôi, hỏi nhỏ: ``\vie{\hl{Mấy ông lại bày trò mèo gì nữa đây?}}''

Tôi nhún vai, cười nhẹ: ``\vie{\hl{Em ấy tự nhiên bảo là muốn làm thám tử. Chắc là xem nhiều phim trinh thám quá nên nổi hứng thôi. Mấy ngày nay tôi có cho em ấy mượn kha khá băng đĩa rồi.}}''

``\vie{\hl{Ý ông là kiểu như \textit{Sherlock Holmes} ấy hả?}}'' Fukuyasu hỏi tiếp.

``\vie{\hl{Ừ, đại loại thế.}}''

Ryuuhou chen vào, mắt vẫn dán vào Shion trong khi nói với chúng tôi: ``\vie{\hl{Em thấy Shion-san nên bỏ phim của anh đi, đọc thêm vài cuốn trinh thám viết bởi các tác giả nổi tiếng còn hơn. Phim của anh cổ lỗ đến mức em nghĩ không ai còn xem nữa đâu, ngoài mẹ với anh ra.}}''

``\vie{\hl{Chính vì nó cổ lỗ nên mới hay chứ em,}}'' tôi xì xào đáp lại. ``\vie{\hl{Mấy cái phim ngày xưa còn hay hơn mớ rác bây giờ.}}''

Towa-sama vẫn không nói gì, nhưng có thể thấy rõ nàng Ác ma đang run nhẹ, có lẽ vì sự ``nguy hiểm'' toát ra từ dáng ngồi ngông nghênh của tiền bối mình. Còn tôi thì cố nhịn cười, nhưng không kiềm được mà bật ra một câu: ``Thám tử đại tài của chúng ta mới chỉ lập nghiệp\dots\ cách đây vài phút.''

Shion ngay lập tức quay ngoắt sang tôi, đôi mắt híp lại đầy nghi hoặc, miệng bĩu ra trông có vẻ khó chịu. Nhưng thay vì phản bác, em ấy chỉ hừ một tiếng đầy bí hiểm, ngó lơ lời của tôi.

Ryuuhou thấy vậy liền khoanh tay, gọi lớn với vẻ mặt tỉnh bơ: ``Ngầu đấy Shion-san, nhưng làm ơn đừng nhìn người khác bằng ánh mắt đó nữa, em sợ chị soi ra tội mà em không hề có đâu.''

Nàng Phù thủy hơi ngượng, nhưng nhanh chóng lấy lại phong độ, chỉ tay vào con bé: ``Còn em, đừng nghĩ tôi không biết em đang quay lén! Mọi ghi hình đều phải xin phép thám tử trước!''

``Rồi, rồi, biết rồi\~{}'' cô em gái chỉ đáp lại đầy suồng sã, chẳng hề mảy may để tâm tới lời của cô nàng.

\ct

Chưa kịp để Shion thể hiện thêm phong thái chuyên nghiệp của mình, một tiếng bước chân vội vã vang lên ngoài hành lang. Pekora, trông hốt hoảng như thể vừa trải qua một cơn ác mộng, lao vào phòng với vẻ mặt đầy tuyệt vọng. ``\textbf{Siêu thám tử! Cứu em với peko!!}''

Tên ``thám tử'' khẽ nhếch môi, nhấc nhẹ chiếc kính một tròng lên rồi đặt nó lại đúng vị trí. Ánh mắt em ấy trở nên sắc lạnh, giọng điệu trầm xuống đầy bí ẩn: ``Sao thế, cừu non của ta?''

Ryuuhou ghé sát vào tai Towa đang đứng cạnh mình, thì thầm: ``\hl{Cừu non? Em tưởng Pekora-san là thỏ cơ mà? Shion-san đang nhầm vụ án nào với vụ nào thế này?}''

Towa nhún vai, cũng thì thầm đáp lại: ``\hl{Chị cũng không rõ nữa, chắc là chị ấy thấy cách gọi này ngầu lòi nên dùng đấy.}''

Tôi khẽ cười, giải thích nhỏ với cả hai cô gái: ``Em ấy thích gọi hậu bối như vậy, vì trong phim hay có kiểu đó. Tin anh đi, em ấy còn bắt Towa-chan gọi mình bằng `đại ca' nữa cơ.''

Towa nghe vậy liền khựng lại, mặt ngơ ngác vì không tin vào tai mình. Tôi thì chỉ tủm tỉm cười, bởi tôi biết chắc rằng, Shion đã ép hậu bối của mình phải gọi như vậy từ trước rồi.

Em tôi nhướng mày, vẻ mặt đầy thích thú: ``Trò này thú vị ghê đấy. Có khi nào em sẽ gọi như thế với các thành viên trong nhóm mới được.'' Tôi cau mày, đổ mồ hôi hột: ``Em định rao giảng đạo giáo đấy à?''

Nàng Thỏ trong khi đó, gần như quỳ xuống sàn, hai tay ôm lấy bím tóc, nước mắt rưng rưng như thể vừa đánh mất vật quý giá nhất đời mình. Giọng em ấy run rẩy: ``Củ cà rốt trên tóc của em biến mất rồi peko!!''

\img{8-11c}

Shion lập tức gật đầu, khoanh tay suy ngẫm trong vài giây rồi phán một câu đầy uy nghiêm: ``OK, cái này thì đúng là vụ lớn rồi.''

Ryuuhou bật cười, lắc đầu nói với Fukuyasu đứng bên cạnh: ``Vụ lớn á? Em tưởng Shion-san sẽ nói gì to tát chứ, củ cà rốt trên tóc mà cũng là vụ lớn à? Mua một củ cà rốt về gắn là được mà.''

Fukuyasu - người vẫn chưa kịp hiểu đầu cua tai nheo ra sao - nhìn quanh, rồi quay sang tôi với vẻ mặt khó hiểu. Tôi chỉ nhún vai, bình thản đáp: ``Ông cứ nhìn tóc của em ấy là đoán được mà.''

Mất một lúc, ông bạn tôi mới vỡ lẽ, gật gù. ``À, cái lọn tóc.''

Trong lúc đó, Shion rút ra một chiếc kính lúp - không ai biết em ấy lấy từ đâu ra - và dõng dạc tuyên bố: ``Bắt đầu điều tra nào!''

Em ấy cúi sát xuống, chăm chú quan sát lỗ trống trên tóc của đàn em với vẻ nghiêm túc quá mức cần thiết. Nhưng từ góc nhìn của Fukuyasu, mọi chuyện lại có vẻ hơi\dots\ khác. ``Sao tôi thấy nó lại ngắm cái gì đó\dots\ dâm đãng thế?''

Ryuuhou nghe vậy gần như phun nước miếng ra vì cười, vỗ vào vai Fukuyasu: ``Trời ơi, Fukuyasu-san nói thế làm em không dám nhìn Shion-san nữa! Sao anh lại nghĩ đến cái chuyện đó được?''

Câu bình luận của cậu bạn khiến tôi khẽ bật cười. Đúng là con mắt sắc bén của một kẻ từng trải.

Sau một hồi ``điều tra'' mà thực ra mắt thường cũng có thể nhận ra ngay, Shion cuối cùng cũng đưa ra kết luận một cách rất tỉnh bơ: ``Chị đã chơi ném phi tiêu với nó ngày hôm qua. Chắc là nó vẫn còn ở đó.''

\img{8-11d}

Towa nghe vậy lập tức trợn mắt: ``Đó là chị sao!? Chị đã làm gì mà để mất nó vậy?''

Ryuuhou cười lớn, vỗ tay khen: ``Thế ra thám tử không chỉ tìm ra thủ phạm, mà còn tự tố cáo luôn mình là người gây ra vụ án đấy. Trình độ điều tra này em thực sự phục lắm!''

Chúng tôi thì chỉ biết thở dài ngao ngán. Cái cách Shion phán như thể đang vạch trần một vụ án động trời, nhưng cuối cùng lại thành tự bóc phốt chính mình. Tôi lắc đầu, than thở với tông mỉa mai: ``Không thể ngờ thủ phạm lại chính là ngài đấy, `siêu thám tử' ơi.''

Pekora không để bụng tới chúng tôi lắm, vui vẻ cảm ơn Shion rồi tung tăng chạy ra khỏi phòng, trên tay vẫn cầm cọng cà rốt vừa tìm thấy. Còn Shion thì ngả người ra ghế, khoanh tay đầy tự mãn, tự hào tuyên bố: ``Không có bí ẩn nào là ta đây không thể giải nổi cả!''

Nhưng sự tự đắc đó nhanh chóng bị dội một gáo nước lạnh khi Towa chép miệng: ``Thế quỷ nào mà chị lại thấy tự hào với cái đó được chứ?''

Ryuuhou nheo mắt, giọng đầy bỡn cợt: ``Shion-san, em nghĩ chị nên thêm một dòng vào danh thiếp của mình: `chuyên giải các vụ án mà chính mình gây ra'. Đảm bảo uy tín oách xà lách luôn.''

Shion quay lại, liếc hai cô nàng vừa khịa mình bằng ánh mắt không thoải mái. Em ấy vẫn còn muốn giữ nguyên phong cách chuyên nghiệp của mình, nhưng bị ngắt quãng thêm một lần nữa bởi nàng Ác ma: ``Và bỏ cái tẩu thuốc ra đi! Hại người lắm đấy!''

Tôi khẽ cười, giải thích nhẹ nhàng: ``Đây là tẩu thuốc đồ chơi thôi mà\dots'' vừa nói tôi vừa giơ tay lên vẫy vẫy, nhưng chẳng ai thèm để ý.

Cô em gái khoanh tay, gật gù đồng tình với nàng Ác ma: ``Mặc dù em thấy cái tẩu đó nhìn cũng khá chất, nhưng Towa-san nói đúng, trông hơi kỳ cục khi cầm cái đó mà mặc đồ phù thủy. Không ăn nhập gì cả.''

Dưới áp lực của cả ba chúng tôi, Shion đành phải miễn cưỡng vứt chiếc tẩu thuốc đi, mà không hề biết rằng thứ mà cô quẳng đi kia sẽ lại kéo cô tới một vụ án khác\dots

Nhưng chưa kịp thở phào, một giọng nói khác lại vang lên ngoài cửa. Lại có người cần sự trợ giúp!

Ngay lập tức, Shion bật dậy, nắm chặt tay, ánh mắt sáng lên đầy quyết tâm: ``Trò chơi đã sẵn sàng kìa! Đi nào, trợ lý của ta!'' rồi hướng mắt về phía chúng tôi, cười tươi: ``Các người cũng đến nào!''

Cả tôi và Fukuyasu đều giơ hai tay lên trời, đồng loạt xua tay như muốn nói ``Bọn anh xin miễn nha.'' Ryuuhou cũng làm động tác y hệt, giọng đầy tinh nghịch: ``Em xin lỗi nha Shion-san, em cũng không có hứng thú với trò chơi thám tử của chị đâu.''

Shion khựng lại một giây, ánh mắt liếc qua chúng tôi, rồi nàng hướng thẳng về phía Towa. Trước khi để đối phương kịp phản ứng, em ấy đã chộp lấy tay của nàng Ác ma và tức khắc kéo đi mất.

\img{8-11e}

Nàng Ác ma giãy nảy: ``Tại sao lại là em chứ!? Tại sao không phải ai khác mà lại là em!?''

Fukuyasu nhún vai, tỉnh bơ đáp: ``Ờ thì, em là người nói chuyện nhiều nhất với `thám tử' mà. Chị em hợp đấy.''

Ryuuhou cũng phụ họa, bước đến gần Towa rồi nói nhỏ giọng trêu chọc: ``Towa-san, tại chị mồm to quá đấy. Lần sau trước khi nói gì phải nghĩ kỹ, rút kinh nghiệm nhé.''

``Nhưng cả ba người cũng nói chuyện với chị ấy mà!?'' Towa lên giọng phản bác, vẫn cố vùng vẫy nhưng vô ích.

Tôi bật cười, vỗ vai Fukuyasu rồi quay sang nói với em ấy: ``Anh không chơi trò này. Với lại, em đang là người mồm to nhất trong phòng đấy, ai ai cũng thấy.''

Ryuuhou cười lớn, rồi nói lớn với Towa đang bị kéo đi: ``Towa-san cố lên! Em sẽ ghi lại toàn bộ buổi điều tra này để làm tư liệu cho buổi hòa nhạc của ban nhạc em! Chắc chắn sẽ rất hay!''

Nghe lời ba người như đang chối bỏ trách nhiệm, Towa chết sững, còn Shion thì kéo nàng đi với khí thế hừng hực. Mặc cho nàng Ác ma cố gắng vùng vẫy thoát, cô cuối cùng vẫn bị kéo đi bởi một ``siêu thám tử'' muốn thể hiện bản thân.

Ryuuhou quay sang tôi, giọng đầy phấn khích: ``Onii-san, em cá là Shion-san sẽ gây thêm ít nhất hai vụ rối nữa trước khi chiều tối. Kiểu gì em phải kiếm thêm hai trợ lý dọn dẹp hậu quả mới xuể.''

Cuộc hành trình phá án của ``thám tử lừng danh'' Shion cùng ``trợ lý bất đắc dĩ'' Towa vẫn sẽ còn tiếp tục.

\ct

Người tiếp theo tìm đến cô nàng ``thám tử'' tự xưng là Choco. So với những người trước, nàng Y tá này lại có một dáng vẻ đầy e thẹn, tay đặt trước ngực, ánh mắt long lanh đầy mong chờ - trái ngược hoàn toàn với sự quyến rũ thườn thường ngày của mình. Giọng nói của cô nhẹ nhàng như gió thoảng: ``Cô rơi mất túi và ví rồi.''

\img{8-11f}

Shion liếc nhìn nàng Y tá bằng nửa con mắt, rồi cũng hạ giọng xuống ngang mức e thẹn như vậy, chậm rãi phán một câu chóng vánh: ``Lập biên bản và báo với cảnh sát đi.''

Căn phòng chìm vào vài giây im lặng. Ryuuhou, vẫn đang ngồi dựa tường ở góc phòng sau khi rời khỏi chỗ ngồi cạnh tôi, bật cười thành tiếng, vừa lắc đầu vừa tiến lại gần bàn làm việc: ``Cái gì thế này? Thám tử mà bảo nạn nhân tự đi gọi cảnh sát á? Vậy Shion-san mở văn phòng này để trang trí hay để làm trò cười cho mọi người thôi vậy?''

Towa sững sờ, miệng há hốc, như thể vừa tận mắt chứng kiến một tội ác không thể tha thứ: ``Liêm sỉ của chị vứt quánh đi đâu rồi ư!?''

Fukuyasu bật cười, khoanh tay lắc đầu: ``Thám tử như này là vứt đi rồi.'' Tôi cũng không nhịn được mà phì cười: ``Thế mà cũng bày đặt làm thám tử nữa!''

Ryuuhou đứng cạnh Choco, vỗ nhẹ vào vai nàng với giọng đầy thương hại, nhưng cũng không giấu được nụ cười nghịch ngợm: ``Choco-san, cô cứ đi báo cảnh sát đi, có khi họ còn hiệu quả hơn đó. Chứ Shion-san mà điều tra thì có khi cô lại mất thêm cả thẻ tín dụng nữa đấy.''

Choco vừa cười nhẹ vừa lắc đầu, vẻ mặt vẫn tươi tỉnh không hề bực bội: ``Cảm ơn em nhắc nhé, cô sẽ cẩn thận.'' Rồi nàng vui vẻ gật đầu với Shion, còn cảm ơn ``thám tử'' trước khi rời đi một cách nhẹ nhàng như thể vừa được ban một lời phán truyền kỳ diệu. Hoặc có lẽ là cô ấy quá dễ dãi. Có lẽ là vì cô biết rằng, ``thám tử'' kia chỉ là một cô bé đang chơi trò chơi của mình. Nếu như vậy thì có lẽ Choco khá là thâm đấy.

Ryuuhou nhìn bóng dáng người rời khỏi phòng, quay sang lắc đầu với Towa, giọng vẫn còn cười: ``Choco-san mà còn kiên nhẫn thế thì em phục thật. Chứ chắc là em thì em đã chửi không kịp vuốt mặt rồi.''

Nàng Ác ma cũng đồng tình, khoanh tay thở dài: ``Đúng là chị ấy có thể làm mọi người mất niềm tin thật\dots''

Shion chắp tay ra sau lưng, thở ra đầy đắc ý: ``Cứ có việc gì đó\dots''

Towa vỗ trán, bất lực: ``\dots chị lại bảo đi nhờ người khác.'' Fukuyasu cũng thêm vào: ``Chưa kể còn đùn đẩy cho người ta nữa, mà lại là cảnh sát nữa chứ!''

Tôi chỉ nhún vai, cười nhạt: ``Anh thà đi ra ngoài tự tìm còn hơn, chứ nhờ mấy bố cảnh sát đó thì đến bao giờ mới thấy?''

Fukuyasu nhướng mày: ``\vie{Cảnh sát ở Nhật vô dụng đến thế á?}''

Tôi cười khẩy, dựa lưng vào ghế: ``\vie{Cứ nhìn cách họ bắt lợn rừng thì đủ hiểu là họ `rác' thế nào rồi đấy. Thua cả mấy ông dân thường ở xứ mình.}''

Ryuuhou, nghe vậy, nhếch mép cười với Fukuyasu: ``\vie{Anh Kuon nói thế tức là em thấy ban nhạc em còn hơn cảnh sát đấy. Bọn em còn thường xuyên tổ chức đội tìm đồ thất lạc cho các thành viên, ít nhất là không bao giờ đẩy người ta đi tự giải quyết vấn đề của mình.}''

Fukuyasu gật gù đồng tình: ``\vie{Cái đó thì đúng thật, có đội ngũ hỗ trợ thì yên tâm hơn nhiều.}''

Towa định phản bác gì đó, nhưng ngay khi em ấy quay sang bàn làm việc, ánh mắt bỗng trợn trừng. Một giây sau, em ấy hét lên trong hoảng loạn: ``\textbf{Có chuyện rồi, Siêu Thám tử ơi!}''

Lập tức, mắt của Shion sáng lên, như thể đã chờ khoảnh khắc này từ lâu. Em ấy giơ tay, hùng hồn tuyên bố: ``Giờ của ta đã đến!''

``\textbf{Thám tử!! Có vụ giết người kìa!!}'' nàng Ác ma hét lên, làm cả phòng giật bắn.

Ryuuhou đứng bật dậy từ vị trí ngồi dựa tường, vẻ mặt hơi kinh ngạc thay vì chỉ có sự ngạc nhiên đơn thuần: ``Vụ giết người? Chị đùa chứ Towa-san?'' Cô vừa nói vừa dời bước về phía bàn làm việc, mắt dừng lại trên thân hình A-chan đang nằm gục bất động.

Chúng tôi lập tức nhìn về hướng mà em ấy chỉ. Và ngay trước mắt, A-san đang nằm gục trên bàn làm việc, bất động, trước mặt là một vũng chất lỏng đỏ rực\dots\ hình như là máu? Không, khoan đã\dots

\img{8-11g}

Tôi liếc nhanh xuống nền nhà. Một lon sốt cà chua bị đổ nằm ngửa ra đó, dòng chữ in trên thân lon vẫn còn rõ ràng giữa vũng sốt đỏ: ``Đến cả ma cũng phải hoàn thành công việc nữa.'' Tôi và Fukuyasu chỉ mất chưa đầy ba giây để nhận ra đây chỉ là một trò đùa. Nhưng Towa thì khác, có lẽ vì quá hoảng sợ nên nàng chẳng hề để ý đến chi tiết nhỏ đó.

Ryuuhou cũng nhanh chóng nắm bắt được tình hình, cô nghiêng người nhẹ nhàng để đọc kỹ dòng chữ trên lon sốt, rồi vỗ nhẹ vào vai Towa, giọng cười vẫn còn run rẩy: ``Bình tĩnh nào, chỉ là trò đùa thôi. Nhìn cái lon này kìa, nó ghi rõ là sốt cà chua mà. A-chan chắc lại bày trò trêu ai đó đấy.''

Tuy nhiên, điều khiến tôi thực sự bất ngờ lại là phản ứng của ``Siêu thám tử''. Người mới đây còn hừng hực khí thế muốn phá án, bỗng dưng đưa tay che mặt, cả người hơi run lên.

Ba chúng tôi nhìn nhau, rồi cùng quay mắt nhìn Shion, hoàn toàn ngơ ngác. Trước khi kẻ nào khác kịp mở miệng, em gái tôi đã bước lại gần cô một bước, tay chạm nhẹ vào cánh tay Shion, giọng vừa ngạc nhiên vừa trêu chọc: ``Ủa, mới nãy còn hùng hồn muốn xử lý vụ án, giờ lại sợ rồi hả thám tử Shion-san?''

\img{8-11g1}

Towa nhíu mày, vẫn chưa hết bối rối: ``Sao thế? Chị vừa nói giờ của chị đã đến mà, mau vào việc đi chứ!'' Shion nuốt nước bọt, cố lùi ra xa một chút, rồi lí nhí nói: ``Chị\dots\ Chị sợ máu lắm, xin lỗi em nha\dots''

Không khí cả phòng lập tức im bặt. Cả bốn người chúng tôi đều không thể tin được vào những gì vừa nghe. Tôi lặng lẽ vỗ tay chậm rãi, đầy mỉa mai: ``Rồi, bó tay hoàn toàn.''

Fukuyasu bật cười lớn, khoanh tay, giọng đầy châm chọc: ``Thậm chí còn thua cả mấy ông cảnh sát nữa thì tôi chả hiểu nổi luôn.''

Ryuuhou ngồi sụt xuống chiếc ghế gần đó, mắt vẫn không rời khỏi Shion với vẻ mặt hoàn toàn không thể tin nổi. Cô quay sang nói với Towa, giọng đủ to để mọi người trong phòng đều nghe thấy: ``Mọi người thấy chưa!? Làm thám tử mà sợ máu thì giải quyết được vụ nào nữa! Trong ban nhạc của em có mấy đứa sợ kim tiêm còn dũng cảm hơn thế này.'' Towa chỉ biết gật gù đồng tình, mặt vẫn còn cứng đờ vì sốc.

Còn tôi thì thở dài, ngửa đầu nhìn lên trần nhà. Đúng là ``Siêu thám tử''\dots\ hết thuốc chữa thật. Ngay sau lời thú nhận hèn nhát đó, nàng Ác ma tức giận, túm lấy hai vai Shion mà lắc tới lắc lui như muốn gọi cô tỉnh lại: ``Chị là kiểu thám tử gì thế này!?''

Cái ``thám tử'' kia lảo đảo, suýt nữa thì ngã sấp mặt, vội vàng hét lên: ``\textbf{Từ từ đã! Chị vẫn còn con bài tẩy nữa mà!}''

Cả bốn chúng tôi đồng loạt nhướng mày. Tôi nhìn Fukuyasu, Fukuyasu nhìn Towa, Towa nhìn Ryuuhou, Ryuuhou nhìn Shion. Cả bốn người đều cùng nghĩ một điều: ``Để xem con bài tẩy này có thực sự ra gì không đây\dots'' Em tôi còn khẽ huých nhẹ vào cánh tay Towa, nháy mắt một cái, như thể hai người đang cùng đặt cược xem Shion sẽ bày ra được trò gì nữa.

Ngay sau đó, Shion giơ tay lên, lẩm bẩm câu thần chú gì đó, rồi một vòng sáng ma thuật lóe lên dưới lòng bàn tay của em ấy. Chỉ trong nháy mắt, từ luồng ánh sáng đó, một sinh vật bốn chân lông xù xuất hiện, tung tăng vẫy đuôi đầy phấn khích.

Đó là một con chó. Một con vật hoàn toàn bình thường, có thể bắt gặp ở bất kỳ ngóc ngách nào trong thành phố Điện tử này.

\begin{figure}[H]
  \centering
  \begin{subfigure}[t]{0.45\textwidth}
    \centering
    \includegraphics[width=8cm]{../../../../Image/Book8/8-11h1.png}
  \end{subfigure}
  \quad
  \begin{subfigure}[t]{0.45\textwidth}
    \centering
    \includegraphics[width=8cm]{../../../../Image/Book8/8-11h2.png}
  \end{subfigure}
\end{figure}

Ryuuhou đứng bật dậy, mắt mở toang vì bất ngờ trước cảnh tượng, rồi bật cười lớn như không kiềm được. Cô tiến lại gần con chó lông xù để vuốt ve, quay sang nói với Shion mà tiếng cười vẫn còn nghẹn ngào: ``Triệu hồi chó ra làm trợ lý thám tử hả? Thật thú vị ghê. Ở ban nhạc của em cũng có một người bạn hay làm y hệt vậy, cứ gặp chuyện khó khăn là lôi mèo cưng của mình ra giải quyết, còn lý luận là 'thú cưng còn thông minh hơn chủ nhân' nữa.''

Shion chống tay lên hông, chỉ vào sinh vật vừa được triệu hồi và tự hào tuyên bố: ``Ta đặt tên cho nó là Thám tử Chó!''

Cả căn phòng lập tức chìm vào im lặng. Bốn người chúng tôi đồng thanh lên tiếng, giọng đầy thất vọng và chán ngán: ``Thế thì chị/em làm thám tử để làm cái gì nữa chứ!?''

Cô nàng Trưởng nhóm vừa nói thêm, vừa dựa lưng vào bàn làm việc, khoanh tay với gương mặt giả vờ nghiêm túc: ``Nói thật nhé, chị có vẻ như đang dựng lên một hệ thống để trốn tránh mọi trách nhiệm đấy. Em thấy chị nên treo biển giải nghệ đi, mở tiệm cà phê mèo còn kiếm được nhiều tiền hơn là làm thám tử.''

Shion phớt lờ tất cả lời châm biếm, vung tay đầy phong thái của một người chỉ huy: ``Nào, Josephine! Bắt đầu công việc thôi!''

Towa mắt tròn tròn lên ngạc nhiên: ``Cái con chó này còn có tên riêng nữa chứ!?''

Fukuyasu khoanh tay, thở dài ngao ngán: ``Không biết con chó này sẽ làm được trò gì nữa đây\dots'' Tôi chỉ cười, lắc đầu đáp lại: ``Đừng có khinh thường nó vội, cứ xem rồi sẽ biết.''

Cô em gái tôi nghiêng đầu nhìn con chó, nhận xét: ``Cái tên Josephine nghe cứ như tên nhân vật nữ chính trong phim cổ trang ấy. Thôi thì cũng coi như Shion-san có gu đặt tên chút.''

Vừa dứt câu, con chó bỗng dưng\dots\ đứng thẳng hoàn toàn bằng hai chân sau. Chưa hết, nó còn lôi từ đâu ra một chiếc kính lúp, chỉnh lại cổ áo (dù nó vốn chẳng mặc gì cả) rồi bước đi điềm nhiên đến chỗ A-chan đang nằm ``gục ngã''.

Cảnh tượng ấy làm tôi nhớ lại con chó đã biến hình sau khi ăn phải một quả cầu kỳ quặc hơn một năm rưỡi về trước\footnote{Xem lại {\flaregothic ÉPISODE 21}, quyển số Hai.}

Ryuuhou mở to mắt vì kinh ngạc: ``Con chó này mà biết đi bằng hai chân á? Ghê quá đi!''

Shion chưa kịp ra lệnh thêm gì, thì trợ lý Josephine đã cúi xuống nhặt lên món đồ mà chủ nhân của nó vừa vứt đi lúc nãy - chiếc tẩu thuốc - rồi quan sát nó một cách tỉ mỉ. Sau vài giây im lặng, con chó hắng giọng, tuyên bố một cách nghiêm túc như thật: ``Cô ấy bị cái tẩu thuốc này đâm trúng đầu.''

\img{8-11i}

Không gian trong phòng bỗng dưng đóng băng.

Không ai cần trao đổi với ai, tất cả chúng tôi đồng loạt ngắt lời, bật ngửa trước cảnh tượng vừa xảy ra. Một con chó vừa bình thản đưa ra phán đoán điều tra, như thể việc phân tích hiện trường là chuyện thường nhật chẳng có gì lạ.

Ryuuhou há hốc miệng, mắt dán chặt vào Josephine trước khi quay sang cậu bạn trên An, giọng còn ngân ngấn sự ngạc nhiên: ``Không thể ngờ là con chó này lại đưa ra phán đoán nhanh như vậy\dots\ Xem chừng nó không phải là con chó bình thường.''

``Em nghĩ có điều gì bình thường ở văn phòng \hll\ này không, Ryuuhou?'' tôi đáp lại với sự bàng quan như đã quá quen với những chuyện vô lý như thế này.

Nhưng những điều bất ngờ còn chưa dừng lại.

Josephine vẫn xoay nhẹ chiếc tẩu thuốc trên bàn chân, rồi tiếp tục với giọng điệu trầm ổn chẳng đổi thay: ``Tôi tin vũ khí này đã được bắn ra từ một khẩu súng, không phải vứt bằng tay.''

Cả văn phòng lại chìm vào im lặng, dày đặc đến mức có thể nghe thấy tiếng hơi thở của mỗi người. Ryuuhou vô thức đưa tay che miệng, đôi mắt mở to đến mức tròn xoe, không thể tin được vào những gì mình vừa chứng kiến. Towa há miệng định nói gì đó, nhưng cổ họng như bị bóp chặt, chẳng thể thốt nên một lời. Fukuyasu đứng khoanh tay, nheo mắt nhìn con chó, vẻ mặt lẫn lộn giữa nghi ngờ và hoang mang.

Còn tôi, tôi thực sự bắt đầu thấy có chút rùng mình trước sinh vật này. Em gái tôi dường như cũng cùng suy nghĩ, cô bước lại gần tôi, thì thầm chỉ đủ hai người nghe: ``Nếu ban nhạc em có một quản lý có khả năng quan sát sắc bén như thế này, em chắc chắn đỡ phải xử lý bao nhiêu chuyện rắc rối mỗi tuần. Nó còn hơn cả cậu quản lý hay quên lịch biểu diễn của bọn em nữa.''

Nhưng điều khiến mọi người bất ngờ nhất, lại là phản ứng của chính người vừa triệu hồi Josephine ra. Shion bỗng chỉ thẳng ngón tay lên trời, hét lớn đến mức kính rơi khỏi mũi: ``\textbf{Một trong số chúng ta ở đây là kẻ ám sát!!}''

Tôi và Fukuyasu đồng loạt thở dài ngao ngán. Ryuuhou tiến lại gần Shion, lắc đầu với ánh mắt vừa thương hại vừa cười nhạo: ``Thế mà chị còn dám nhận mình là thám tử chính? Josephine mới là người vừa làm hết tất cả công việc phân tích hiện trường đấy. Chị định ngồi đây hưởng thành quả của trợ lý à?''

Tôi đưa tay xoa trán, cố gắng giữ bình tĩnh trước sự hời hợt của cô nàng phù thủy: ``Bọn anh biết ai là thủ phạm rồi.''

Shion lập tức mở to mắt, giọng gấp gáp như sắp nhảy lên: ``Là ai!? Nói nhanh đi!'' Tôi không vội trả lời, chỉ khẽ mỉm cười khẩy: ``Cứ chờ cảnh sát đến đi, rồi em sẽ tự hiểu.''

Shion nheo mắt đầy nghi ngờ, lùi lại một bước: ``Cảnh sát nào cơ? Anh gọi ai thế?''

Ngay sau câu hỏi của cô, Fukuyasu lặng lẽ giơ tay, chỉ về phía góc phòng nơi Towa đang đứng.

Ở đó, ``trợ lý'' của Shion - chính là Towa, người vẫn giữ được sự tỉnh táo xuyên suốt toàn bộ vụ việc - đang bình thản cầm điện thoại, bấm số gọi người đến giải quyết chuyện này.

\img{8-11j}

Ryuuhou bật cười thành tiếng, tiến lại vỗ nhẹ vai Towa trước khi quay sang Shion: ``Nói mới nhớ, trợ lý đắc lực nhất của chị đây này. Towa-san từ nãy đến giờ mới là người xử lý mọi việc đúng nghĩa, Josephine đâu phải chỉ mình nó có ích đâu.''

Không cần phải nói ra, ai ở đây cũng rõ ràng thủ phạm là ai - trừ cái tên ``thám tử'' dởm và ngây thơ kia ra.
%==========================================TRUYỆN PHỤ=========================================
\omk

Ngay sau cuộc gọi của Towa, cánh cửa văn phòng bật tung ra, một nhóm ``cảnh sát'' nhanh chóng chiếm lĩnh toàn bộ không gian. Một giọng nói quen thuộc vang lên, đầy khí thế: ``Ở đây có ai báo án mạng hả!?''

Tôi quay đầu lại theo phản xạ, và lập tức trố mắt khi nhìn thấy người đứng trước cửa. Nếu đó là một sĩ quan cảnh sát chính quy, tôi cũng chẳng ngạc nhiên đâu. Nhưng đằng này, người vừa hùng hổ bước vào, mặc ``đồng phục'' gồm bộ đồ thường ngày kèm chiếc mũ cảnh sát đội chễm chệ trên đầu\dots\ lại là Subaru.

Đúng vậy. Nàng Vịt của chúng ta đã chính thức-- à không, tạm thời ``nhậm chức'' cảnh sát dạo rồi.

Ryuuhou từ góc phòng nhướng mày quan sát Subaru, rồi cười khẽ với cô nàng đang đứng cạnh mình là Towa: ``Chị ấy diễn nhập tâm quá. Trông hợp làm sao ấy.''

Tất nhiên, đây chỉ là một trò chơi trong văn phòng. Nhưng nếu phải chọn giữa gọi cảnh sát thật hay gọi Subaru, tôi thà chọn em ấy còn hơn. Lý do? Cảnh sát thật mà đem ra xử lý mấy trò đùa này, có khi còn kém cỏi hơn nàng Vịt này nhiều.

Vừa thấy Subaru bước vào, Fukuyasu vỗ vai tôi, cười cợt: ``Xem ra lực lượng tinh nhuệ đã đến nơi rồi.'' Towa thì khoanh tay, gật gù đầy tâm đắc: ``Cảm giác này đúng là an toàn hơn hẳn.''

Cô em gái của tôi đứng dậy, tựa tay lên bàn, nhận xét với giọng trêu chọc nhưng có vẻ suy nghĩ nghiêm túc: ``Subaru-san mà làm cảnh sát thật, chắc không tội phạm nào dám chống đối đâu. Cái chất giọng đó vừa thấy đã muốn đầu hàng rồi.''

Riêng Shion thì mặt bắt đầu tái đi, có một dự cảm chẳng lành lóe lên: ``Ơ\dots\ đợi đã\dots''

Nhưng chẳng kịp để em ấy phản ứng gì thêm, tất cả chúng tôi đồng loạt chỉ tay về phía nàng Phù thủy, cùng với lời ``khai báo'' chắc chắn từ con chó trợ lý thám tử của chính em ấy. Subaru nhìn Shion, rồi quay sang nhìn Josephine. Con chó cũng nhìn Subaru, rồi lại ngẩng lên nhìn chủ nhân của nó với vẻ mặt thản nhiên. Shion thì nhìn chúng tôi, mặt méo xệch chẳng thể nói nên lời.

Và rồi\dots

*CẠCH!*

Không cần nói thêm lời nào, Subaru ``tặng'' ngay cho Shion một còng số tám sáng loáng, liền kéo đi mất.

\img{8-11k}

``\textbf{K-Khoan dã!!}'' nàng ``Thám tử'' dởm hoảng loạn, cố gắng giãy giụa trong khi bị kéo đi. ``Tôi có làm gì sai đâu chứ!?''

Subaru lắc đầu, đóng vai chuyên nghiệp hết cỡ: ``Mọi người ở đây đều khai cô là thủ phạm mà!''

``Không có! Chỉ là tôi vô tình thôi!''

``Miễn bình luận. Về đồn rồi sẽ nói chuyện kỹ càng. Đi theo tôi!''

Ryuuhou khoanh tay đứng nhìn cảnh tượng, vẻ mặt hài lòng vô cùng. Cô quay sang nói với Fukuyasu đang đứng cạnh: ``Cảnh tượng này đáng để ghi hình lại làm kỷ niệm quá. Thám tử tự phong bị bắt bởi cảnh sát Vịt - có khi em viết thành kịch bản cho buổi hòa nhạc tới của ban nhạc cũng hay nhỉ?''

``Sao lại nói mấy cái này cho anh?'' cậu bạn tên An vừa cười gượng với cô em gái Trưởng nhóm, vừa cảm thấy hả dạ khi tên ``tội phạm'' được dẫn đi.

Shion chỉ có thể gào thét trong vô vọng, cố vùng vẫy nhưng không tài nào thắng nổi sức mạnh của Subaru. Em ấy còn quay lại nhìn chúng tôi với ánh mắt tuyệt vọng pha lẫn phẫn nộ, gào lên hết cỡ: ``\textbf{CÁC NGƯỜI LÀ LŨ BÁN ĐỒNG ĐỘI!!}''

Lời kêu cứu hoàn toàn bị phớt lờ. Ryuuhou phì cười, rồi quay sang tôi, mắt nhính lại đầy tinh nghịch: ``Câu nói ấy hay quá.''

Tiếng gào thét của Shion càng lúc càng nhỏ dần, cho đến khi hoàn toàn biến mất sau cánh cửa văn phòng. Trước khi bị kéo đi khuất hẳn, em ấy còn cố ngoái đầu lại mắng con chó Josephine, đứa vừa phút trước còn là trợ lý trung thành: ``\textbf{Josephine, mày là đồ phản bội!!}''

Con chó chỉ nhún vai đáp lại một cách thản nhiên: ``Công lý là công lý thôi, thưa cô.''

Ryuuhou nhìn cảnh đó, gật gù tán thưởng với Josephine: ``Con chó này có phẩm chất hơn hẳn chủ nhân nó. Có khi em nên mượn nó về làm quản lý cho ban nhạc, vừa có mắt quan sát tốt lại vừa không bao giờ thiên vị ai.''

Đúng là chó cảnh sát có khác, phẩm chất thật sự khác biệt.

Trong khi đó, Towa vẫn ôm lấy A-chan - người vẫn còn bất tỉnh trên ghế - và chào theo kiểu nhà binh khi nhìn bóng dáng Subaru kéo Shion đi. Còn tôi thì chỉ vẫy tay, cười mếu: ``Tí nữa bọn anh lên đồn đón em về!''

Shion vẫn cố gào lên cái gì đó, nhưng tiếng nói ấy càng lúc càng xa, rồi dần dần biến mất hoàn toàn.

Ryuuhou quay sang tôi, giọng đầy hài hước nhưng cũng có chút suy nghĩ: ``Anh biết không, em cá cược là Shion-san sẽ mất ít nhất ba ngày mới dám ngẩng mặt nhìn chúng ta. Sau vụ này thì chị ấy còn chẳng biết mặt mũi để đâu luôn!''

Tôi thở dài, tự hỏi với chính mình: ``Không biết có lần nào mà Shion-chan làm thám tử mà không bị cảnh sát bắt đi không nhỉ?''

Con bé nhún vai, cười khẩy: ``Em nghĩ là không có đâu. Nhưng mà em thích trò này đấy.''

Nhưng\dots\ mọi chuyện vẫn chưa dừng lại ở đó. Ngay khi ``thám tử'' dởm vừa bị lôi đi khuất dạng, Towa lập tức thở phào nhẹ nhõm, nói với giọng khoái chí: ``Cuối cùng cũng tống khứ được tên thám tử vô dụng kia đi được rồi.''

Tôi lập tức quay sang nhìn em ấy, hoàn toàn sững sờ. Ryuuhou cũng dừng lại, mắt mở to bất ngờ, cô bước tới gần cô nàng: ``Khoan\dots\ ý chị là\dots\ chị biết từ đầu đây chỉ là một trò đùa à? Không phải chị cũng bị dính vào tình thế đó sao?''

Nàng Ác ma đáp lại một cách tỉnh bơ, thản nhiên: ``Tất nhiên. Đó cũng là một phần trong trò chơi khăm của chị mà.''

Fukuyasu đứng ngay bên cạnh cũng trợn mắt, không thể tin vào tai mình: ``Khoan, vậy cái dòng chữ ghi sốt cà chua dưới bàn đó là chị sắp đặt hết?''

Towa mỉm cười đầy ẩn ý: ``À, cái đó thì\dots''

Ryuuhou nhìn cô nàng với vẻ mặt vừa thán phục vừa cười ngất, cô vỗ vào vai nàng ác ma: ``Towa-san giấu kỹ quá. Em tưởng chị là nạn nhân của vụ đùa này, hóa ra chị là đạo diễn chính đằng sau. Thế, đầu đuôi mọi chuyện là thế nào vậy?''

Chúng tôi sắp được nghe một màn giải thích hết sức quái thai về toàn bộ kế hoạch.

\begin{center}
  \uline{Hồi tưởng.}
\end{center}

Theo câu chuyện mà Towa kể lại, toàn bộ sự việc bắt đầu từ cái khoảnh khắc Shion vẫn còn đang ngất ngây trong cái chiến thắng giả tạo của bản thân, mải miết chỉnh sửa chiếc mũ và kính lúp như thể mình là một thám tử chính hiệu trong những bộ phim trinh thám.

Còn có một người khác thì sao?

Đó là A-san.

Chị ấy vẫn đang chăm chỉ cắm đầu vào công việc, hoàn toàn làm ngơ trước những trò hề mà ``Thám tử'' đang bày ra. Cặp mắt của chị ấy dán chặt vào màn hình máy tính, ngón tay gõ phím đều đặn không ngừng, như thể toàn bộ thế giới xung quanh chẳng hề tồn tại.

Nhưng đúng lúc ấy, mọi thứ bắt đầu thay đổi. ``A-chan, em có thể nhờ chị làm giúp một việc được không?''

Chị đồng nghiệp ngẩng đầu lên, thoáng chút bối rối: ``Có chuyện gì vậy, Tokoyami-san?''

Towa lập tức giơ lên một lon đồ uống, nét mặt ngây thơ cười nói: ``Em mang theo một lon sốt cà chua đây. Chị có muốn uống bồi bổ lại sức không?''

\dots Cả căn phòng chìm vào im lặng.

Chị ấy nhìn chằm chằm vào thứ nàng Ác ma đang cầm, như thể đang cân nhắc giữa việc nhận lời hay bảo em ấy cất đi. Và rồi, với một nhịp độ cực kỳ chậm chạp\dots

``\dots Chị xin nhé.''

Towa trợn tròn mắt ngạc nhiên. Cô tự nhủ thầm: ``\textit{Chị ấy thực sự đồng ý uống lon sốt cà chua này á!?}'' nhưng chẳng mấy chốc cô đã lấy lại bình tĩnh, nghĩ trong lòng: ``\textit{Chắc chị ấy mệt mỏi quá mức rồi\dots}'' và không để lỡ cơ hội quý giá này, một kế hoạch hoàn hảo lập tức lóe lên trong đầu cô.

Một kế hoạch mà không ai có thể ngờ tới.

Cô đã lén mở lon sốt cà chua, sau đó cố tình làm đổ nó xuống mặt bàn. Thế nhưng, có vẻ như A-chan đã kiệt sức đến mức chẳng thể nhận ra chuyện gì vừa xảy ra, vẫn tiếp tục miệt mài làm việc.

Fukuyasu nhìn Towa, giọng cắt ngang câu chuyện đầy bất ngờ: ``Vậy là chị ấy hoàn toàn không phát hiện ra vũng chất lỏng màu đỏ đó?''

Towa cười một cách bí ẩn, đáp lại: ``A-chan còn mải làm việc nhiều hơn cả em tưởng tượng. Chị ấy cứ như một cỗ robot không biết mệt mỏi vậy.''

``Giống như ai đó trong lớp mình ghê\dots'' cậu bạn và cả em gái tôi đều liếc về phía tôi, cũng là một nạn nhân (tự nguyện) của tư bản. Tôi chẳng cần phải phản biện gì, vì đúng là như vậy; tôi đã không dưới vài ba lần thức xuyên đêm chỉ để hoàn thành công việc, nhưng ít nhất vì vẫn may mắn hơn A-san vì được trả lương thêm giờ, trong khi chị ấy thì hoàn toàn làm không công. Tôi chẳng hiểu vì sao Tanigo-san lại đãi ngộ hai chúng tôi theo hai cách khác nhau như vậy.

Quay trở lại với câu chuyện. Nhận thấy cơ hội quá thuận lợi, Towa nhanh chóng lên tiếng: ``A-chan, chị có muốn viết lại vài dòng tâm sự gì đó không?''

A-chan với vẻ mệt mỏi đáp: ``Ừm\dots\ được thôi.''

Và thế là, trong khi chủ nhân của bàn tay chẳng còn đủ tỉnh táo, Towa đã sắp đặt hết tất cả các ``bằng chứng'' ngay trên mặt bàn, dùng chính tay của A-chan để viết ra dòng chữ đó. Nghe đến đây, Ryuuhou nhìn Towa với vẻ mặt vừa mỉa mai vừa khâm phục: ``Em không hiểu Towa-san đã học được những tuyệt chiêu này từ đâu, nhưng mà chị cũng lắm mưu đấy. Lừa được tất cả mọi người ở đây.''

Thật kỳ lạ làm sao, cả Shion và con chó Josephine đều chẳng hề để ý đến vũng ``máu giả'' này ngay từ đầu. Cả hai chỉ mải mê chạy theo cuộc điều tra hão huyền mà chẳng bao giờ chịu liếc mắt nhìn những chi tiết hiển nhiên ngay trước mặt.

Và rồi, một ``sự cố tình cờ'' đã xảy ra.

Khi Shion ném chiếc tẩu thuốc đi theo yêu cầu của Towa, nó đã trúng ngay vào đầu A-chan, khiến chị ấy gục xuống bàn bất tỉnh tại chỗ. Và thế là\dots\ đầu của chị ấy vô tình rơi ngay sát vũng sốt cà chua. Và, tada! Một vụ ``án mạng'' hoàn hảo đã ra đời!

\begin{center}
  \uline{Kết thúc hồi tưởng.}
\end{center}

Khoan đã, điều này có thật không?

Tôi nhìn Towa, chẳng thể tin được vào những gì mình vừa nghe: ``Em đang nói với bọn anh là\dots\ em đã giăng bẫy tất cả từ đầu?''

Towa mỉm cười đầy tự hào, đưa tay vén nhẹ một lọn tóc: ``Đúng á!''

Tôi lập tức quay sang nhìn Fukuyasu, gương mặt cậu bạn cũng ngẩn ngơ chẳng kém gì tôi. Còn Ryuuhou thì đã cười đến mức chảy nước mắt, vừa đập tay xuống đùi vừa lau nước mắt: ``Trời ơi, Towa-san đúng là một kẻ quái thai. Nếu em mà là Shion-san, em thật sự chẳng còn mặt mũi nào để nhìn mọi người nữa.''

Tôi lắc đầu, thở dài bất lực: ``Towa-chan đúng là quái thai thật đấy.'' Đáp lại lời tôi chỉ là một nụ cười ngập tràn tự hào từ nàng Ác ma.

Ryuuhou tiếp tục nhìn nàng Ác ma, giọng nói vừa mỉa mai vừa đầy sự tán thưởng: ``Towa-san, em thật sự nể phục chị. Mà này, lần sau nếu ban nhạc em có cần dựng một vở kịch nào đó, chị có thể làm đạo diễn cho bọn em không? Em trả công bằng bánh kẹp thịt đầy ụ nhé.''

Chắc chắn nếu Shion mà biết được toàn bộ sự thật, nàng sẽ ``giết chết'' Towa cho mà xem.

\ct

Dù sao đi nữa, cái vụ án mạng giả mạo kia cũng chính thức ``đóng hộp'' rồi. Tôi chống cằm, đang định tổng kết tình hình thì chợt nhớ ra một điểm chưa rõ, liền hỏi Towa: ``Mà này, lúc nãy em đã gọi cho ai đến ``xử lý'' vụ án thế? Làm sao mà đến nơi nhanh thế?''

Nàng Ác ma vẫn bình thản như không có chuyện gì, đáp gọn lỏn: ``Em gọi Subaru-senpai thôi mà.''

Tôi lập tức méo mặt, mắt mở to hết cỡ vì bất ngờ, chưa kịp mở miệng thì Ryuuhou đã bật cười lớn từ bên cạnh, vừa vỗ vào vai tôi vừa nói: ``Trời ơi, chị tính trước hết cả chứ!? Biết chắc là chị Vịt sẽ đến xử lý tận nơi luôn hả?''

Towa nhún vai, giọng vẫn thản nhiên như đang bàn luận thời tiết hôm nay: ``Thì chị ấy là cảnh sát dạo nổi tiếng mà, có vấn đề gì gọi là có liền, đâu phải ai cũng được như vậy. Đáng tin cậy lắm chứ!''

Tôi vẫn chưa hết sốc, lắc đầu không hiểu nổi cặp đôi này đã lên kế hoạch từ lúc nào. Em tôi cũng gật gù đồng tình với tôi, nghiêng người nói nhỏ: ``Chị này thật sự có óc tổ chức thật đấy.''

Đúng là không sai gì cả, nhưng mà tôi vẫn đang thắc mắc làm sao mà Subaru lại có thể mặc bộ ``đồng phục'' cảnh sát đến sớm thế nhỉ? Tôi vẫn nghĩ đó chỉ là một bộ đồ chơi đóng kịch mà thôi, không ai ngờ nàng lại mang đến dùng thật trong tình huống này.

Towa lại thêm vào, giọng đầy tự tin: ``Tí nữa bọn mình lên `đồn' đón chị ấy về là xong. Kuon-senpai yên tâm đi, em nghĩ Subaru-senpai cũng không làm gì quá đáng với Shion-senpai đâu, chỉ đùa một chút thôi mà.''

Lúc này Fukuyasu mới nhíu mày, chỉ về phía A-san vẫn đang ngất xỉu lặng lẽ trên ghế, thắc mắc: ``Còn A-chan thì sao bây giờ? Chị ấy vẫn chưa tỉnh lại nữa.''

Towa chỉ ngay về phía hành lang phía cuối phòng, giọng rành rọt không chần chừ: ``Công ty mình có phòng y tế mà, có giường nghỉ và thuốc men đầy đủ. Ta dìu chị ấy xuống đó nghỉ ngơi đi là ổn thôi.''

Ryuuhou lập tức giơ tay tình nguyện, vừa đi tới phía ghế của chị ấy vừa nói: ``Đúng đó, em và Towa-san sẽ dìu chị ấy xuống, hai người dọn dẹp chỗ này một chút đi, vũng sốt cà chua kia chưa lau kìa.''

Tôi thở dài, đưa tay xoa thái dương mệt mỏi. Thôi được rồi, cứ tạm gác mọi chuyện qua một bên đã. Nhưng có một điều tôi ghi nhớ mãi trong đầu, từ nay về sau, tuyệt đối sẽ không bao giờ tin vào bất kỳ ``vụ án'' nào mà Towa dàn dựng lên nữa. Con bé này có óc sắp xếp quá, lần sau bị nó đem ra làm nhân vật chính trong trò khăm nào cũng không biết.

\ct

Chúng tôi cùng nhau dìu A-san xuống phòng y tế, đắp chăn nhẹ để chị ấy nghỉ ngơi; căn phòng chỉ còn tiếng quạt trần đều đều và nhịp thở thản nhiên của người nằm. Vụ rắc rối vừa rồi dù sao cũng qua đi, và tôi sẽ xử lý nốt công việc còn dang dở của A-san, để chị ấy được yên tâm nghỉ ngơi. Trong lúc bước ra khỏi phòng, Ryuuhou vô tình chạm vai Fukuyasu, vừa cười khúc khích vừa chỉ vào vết sốt cà chua còn vương trên cổ áo cậu bạn: ``Fukuyasu-san, dọn đi kẻo người khác hiểu nhầm là bị thương gì đó đấy. Chị A tỉnh lại mà thấy cảnh này, chắc là chị sẽ quy tiên vì đột quỵ đấy.'' 

Fukuyasu chỉ biết xoa đầu cười trừ, vội vàng lấy khăn giấy lau đi vết bẩn.

Vừa bước ra khỏi phòng y tế, tôi chợt nhớ ra con chó thám tử Josephine vẫn còn ở văn phòng, bèn quay lại để làm rõ một điều khó hiểu trong đầu. Cuộc trò chuyện bất ngờ hé lộ rằng con chó này đâu chỉ là thú cưng thông thường; nó thực sự là một nhà điều tra tự do, Shion đã thuê nó làm trợ lý riêng cho các vụ án trinh thám rườm rà của mình suốt mấy tháng qua.

Cậu bạn là người đầu tiên đặt câu hỏi, khoanh tay trước ngực với giọng thắc mắc không giấu được: ``Vậy mà cậu không nhận ra thứ chất lỏng đỏ trên bàn hôm nay là sốt cà chua à? Cái lon còn để ngay đó cơ mà.''

Nghe câu ấy, Josephine giật nảy mình, đôi tai dựng đứng hẳn lên, giọng đầy kinh ngạc như vừa bị lừa một ván lớn: ``Khoan đã, chuyện đó là thật á!?''

Tôi nhún vai, chỉ vào cái lon rỗng còn nằm dưới gầm bàn tôi vừa mang theo: ``Ừ, ghi rõ chữ 'sốt cà chua' trên nắp lon mà. Cả chúng tôi đều thấy từ đầu.''

Con chó chớp mắt liên tục, cụp đôi tai xuống đầy xấu hổ, nghiêng đầu bối rối hỏi: ``Tôi\dots\ tôi không để ý chút nào. Vậy vừa rồi tôi đã làm sai gì à?''

Tôi lắc đầu, một nụ cười nhẹ nở trên môi: ``\dots Không hẳn. Ngược lại nữa là đằng khác.'' Rồi tôi bật cười nhỏ, thêm vào: ``Cậu làm tốt hơn nhiều người ở đây đấy.'' Fukuyasu cũng gật gù đồng tình, vỗ vai con chó đầy tán thưởng: ``Đúng thế. Thật sự cậu còn có tố chất thám tử hơn cái cô chủ tóc xám của cậu gấp trăm lần.''

``Thế á?'' Josephine ngẩn người ra, có vẻ đây là lần đầu nó được khen một cách chân thành như vậy. Cái đuôi vốn cụp xuống bỗng vẫy nhẹ, dần dần nhanh hơn hẳn.

Tôi tiếp tục khuyến khích: ``Chỉ cần rèn luyện thêm khả năng quan sát chi tiết, cậu hoàn toàn có thể trở thành một thám tử giỏi thực thụ. Không phải kiểu chơi đùa như Shion-chan đâu.''

Lúc này Ryuuhou mới bước lại, dựa lưng vào bàn làm việc với nụ cười nghịch ngợm, nhìn con chó đầy thích thú. Em không chỉ nói với tôi mà còn quay sang giỡn với Fukuyasu trước khi nói với Josephine: ``Tôi cũng ủng hộ ý đó. Cậu này có khí chất hơn Shion-san nhiều, từ lúc phân tích hiện trường đến giọng điệu đều ra dáng nhà điều tra lắm. Này cậu biết chơi nhạc cụ nào không? Ban nhạc của tôi đang thiếu một người quản lý có mắt tinh đời thế này, lương bằng thịt nướng đủ loại, có muốn suy nghĩ không?''

Fukuyasu bật cười lớn, chỉ vào Ryuuhou: ``Em còn định mướn cả chó làm quản lý nữa hả? Thật sự là không bỏ lỡ cơ hội nào.'' Cô bé chỉ nhún vai cười trừ, vẻ mặt rõ ràng là đang nói đùa nhưng cũng có chút thật lòng.

Josephine im lặng một lúc, ngẫm nghĩ những lời chúng tôi vừa nói, vẻ mặt dần trở nên nghiêm túc. Bất ngờ nó cúi đầu cảm ơn rối rít, chào tạm biệt tất cả rồi tức tốc chạy ra khỏi văn phòng, bước chân dồn dập như vừa tìm thấy mục đích sống mới.

Tôi và Fukuyasu đứng nhìn nhau, cả hai đều không hiểu chuyện gì vừa xảy ra. Ryuuhou thì ôm bụng cười lớn, nhìn theo bóng con chó biến mất ở cuối hành lang: ``Thề có chúa, có khi tuần sau nó thực sự mở văn phòng thám tử riêng cạnh công ty mình đấy. Tôi mà thấy biển hiệu `Thám tử Josephine' là đầu tiên đến đăng ký làm khách hàng đầu tiên.''

Cứ thế, một chuyện kỳ lạ nữa vừa trôi qua trong văn phòng của chúng tôi. Nhưng vẫn còn một nhiệm vụ quan trọng đang chờ đợi: phải đến đồn cảnh sát của Subaru, đón cô ``thám tử dởm'' Shion về nhà.

Khi chúng tôi bước qua cửa phòng chờ, cảnh tượng đầu tiên đập vào mắt tôi là\dots

Shion đang ngồi ỉu xìu trong góc phòng, lưng dựa vào tường, hai tay ôm đầu gối. Ngay khi thấy bóng dáng tôi, em ấy lập tức nhổm dậy, đôi mắt sáng rực như vừa thấy vị cứu tinh, rồi lao tới bám lấy cổ áo tôi: ``Cuối cùng anh cũng đến! Tại sao không cứu em sớm hơn!?'' - em ấy ưỡn mặt ra, giọng đầy oan khuất như một đứa trẻ bị bỏ quên. Chưa kịp tôi mở lời, em ấy đã đổi giọng cầu cứu, nghẹn ngào nắm chặt tay áo tôi: ``\textbf{Làm ơn, đừng bỏ em lại đây một mình điiiii!}''

Tôi chỉ thở dài cười nhạt với phản ứng quá đà của em ấy, đúng là em ấy bị oan thật, nhưng cái oan đó hoàn toàn là tự tìm đến. Ryuuhou đứng cạnh tôi, tay khoanh trước ngực, mỉm cười hơi mỉa mai nhưng đầy tinh nghịch: ``Nếu chị cứ tiếp tục theo nghiệp thám tử kiểu này, e rằng phải sắm sẵn bộ đồ đi tù thôi Shion-san ạ. Đừng nói là phòng tạm giam, đến nhà tù lớn chị cũng sẽ còn quen mặt nữa.''

Subaru vừa từ trong phòng giấy bước ra, tay vỗ đùi cười lớn khi nghe câu đó: ``Nghe đâu rồi đấy! Cảnh sát nhà này xử lý tội phạm lắm, không bao giờ tha thứ cho kẻ xấu!'' Ryuuhou liền vỗ vai nàng Vịt, giọng khen ngợi hết mực: ``Subaru-san làm cảnh sát quá đỉnh. Nếu có bảng xếp hạng cảnh sát dạo tốt nhất, chắc chắn chị đứng nhất nhì bảng.'' Subaru ngẩng cao đầu đầy tự hào, còn tôi chỉ lắc đầu nhìn hai người đùa cợt.

Sau một hồi kỳ kèo, thương lượng, cuối cùng Shion cũng được ``tha bổng'' sau khi Subaru hài lòng với màn diễn của mình. Tôi thậm chí không phải tốn một đồng nào để chuộc em ấy ra, chỉ mất mấy lon nước ngọt cho cô nàng cảnh sát dạo là xong.

Nàng Phù thủy vẫn còn đứng bực mình ở cửa ra vào, lẩm bẩm trong miệng, vẻ mặt hậm hực chẳng chịu thua: ``Rõ ràng em là một thám tử lừng danh lắm mà\dots\ Chỉ là không may bị mắc bẫy thôi.''

Fukuyasu vỗ vai em ấy, giọng an ủi nhưng vẫn xen lẫn trêu chọc: ``Trở thành thám tử không dễ như em tưởng đâu, Shion-chan à. Cái nghiệp đó đòi hỏi nhiều hơn là chỉ đội cái mũ và kính lúp đấy.''

Chúng tôi vừa đi dọc hành lang về văn phòng, vừa dành thời gian giải thích cho em ấy cách một nhà điều tra thực thụ vận hành ra sao. Nhiều lúc họ còn phải đánh đổi cả tính mạng chỉ để có được thông tin quan trọng từ đối tượng. Tôi kết luận câu chuyện, giọng nghiêm túc hơn một chút: ``Nói tóm lại, làm thám tử là cực kỳ khó, phải có tinh thần thép mới có thể trụ được lâu dài.''

Ryuuhou bước cạnh Shion, tay đút túi quần, gật gù thêm vào: ``Em nghĩ chị nên bắt đầu từ những bài tập nhỏ trước đã. Ví dụ như tuần nào cũng dành một buổi tìm cái chìa khóa hay điện thoại hay bị lạc của chị ấy. Làm được đều đặn nửa năm, rồi hãy nghĩ đến việc điều tra vụ án lớn. Từ từ mới lên cấp được.''

Sau một đoạn im lặng ngượng nghịu, Shion hắng giọng, rồi quay sang tất cả chúng tôi, lí nhí xin lỗi, đôi má đã hơi ửng đỏ vì xấu hổ khi nhận ra mình đã làm hỏng hết cả kế hoạch.

Tuy nhiên, tôi chỉ mỉm cười và lắc đầu, chỉ về phía Towa đang đi cuối đoàn: ``Anh nghĩ người cần xin lỗi ở đây phải là Towa-chan mới đúng chứ?''

Ngay lập tức, Towa khựng lại bước chân, mặt tái đi khi bị điểm danh.

``Khoan, gì cơ!?'' Shion thảng thốt, mắt mở to hết cỡ, chưa kịp nối được mạch chuyện.

Tôi nhướng mày, nhún vai tỉnh bơ đáp: ``Nếu không có trò đùa quái thai của em ấy dàn dựng từ đầu, thì làm gì có chuyện em bị bắt nhầm làm thủ phạm chứ.''

Shion chớp mắt mấy lần, rồi tròn mắt nhận ra vấn đề, lập tức quay sang nàng Ác ma, giọng đầy phẫn nộ: ``\textbf{Là lỗi của em! Mau xin lỗi chị điiii! Em dám hãm hại cả thám tử của mình à!?}''

Towa nhăn mặt, biết mình không thể chối cãi, đành chịu thua, giọng lí nhí như muỗi kêu: ``\dots Xin lỗi chị nha, Shion-senpai\dots Lần sau em sẽ không để chị bị bắt lâu nữa.''

Nghe câu xin lỗi còn chưa hết tình, Ryuuhou bật cười lớn, bắt gặp ánh mắt của Fukuyasu cũng đang cười khúc khích bên cạnh, rồi lắc đầu nhận xét: ``Đúng là một vòng luẩn quẩn không thể hiểu nổi. Thám tử tự phong bị chính cộng sự của mình dàn cảnh hãm hại, rồi lại bị cảnh sát Vịt bắt giam, cuối cùng phải đợi đoàn bạn bè đến chuộc về. Chuyện này mà em viết thành kịch bản cho buổi hòa nhạc tới của ban nhạc, chắc chắn bán hết vé sớm một tuần. Em còn định mời cả Subaru-san đến diễn chính vai cảnh sát nữa.''

Subaru ở đầu đoàn nghe vậy liền giơ tay tán thành: ``Đồng ý! Tôi diễn rất hay!'' Tôi chỉ có thể đứng giữa mọi người, vừa ngăn cản cặp đôi kia không quá đà trêu chọc Shion, vừa giúp hai cô nàn làm hòa. Cuối cùng, mọi chuyện cũng được giải quyết ổn thỏa, mọi người tiếp tục bước đi trên hành lang dài về văn phòng, tiếng cười nói vang vọng khắp nơi.

\il

Cuối cùng mọi thứ cũng quay về guồng quay cũ, những ồn ào vừa rồi lắng xuống hết như bọt biển tan trên bờ cát. Tiếng cười, tiếng la hét, tiếng gọi nhau í ới trong phòng giờ chỉ còn là âm vọng mờ nhạt. Tôi đứng dậy, vươn vai vài cái, cảm nhận từng khớp xương kêu răng rắc sau một ngày dài rối ren.

Nhưng chẳng hiểu sao, trong lòng tôi dậy lên một cảm giác rất lạ, khó nói thành lời. Nó cứ trôi lơ lửng như bụi trong ánh nắng chiều, hiện hữu mà không thể nào nắm bắt.

Không phải là sự nhẹ nhõm khi một vụ rắc rối vừa được giải quyết xong xuôi, cũng chẳng phải sự mệt mỏi tích lũy sau một ngày dài đầy chuyện phiền phức. Mà nó giống như\dots\ có một sự thay đổi lớn sắp ập đến. Một nỗi lo sợ về một sự mất mát nào đó đang lơ lửng, dần dần đến gần như một cơn bão mà ta có thể cảm nhận qua độ ẩm của không khí, nhưng chưa thấy bóng dáng của nó trên chân trời.

Đang khi tôi miên man suy nghĩ, cơn gió ngoài trời bất chợt nổi lên mạnh hơn hẳn. Nó gào thét qua những khe cửa, mang theo hơi ấm của một mùa hè đang đến gần. Những tấm áp phích dán trên bảng thông báo trước cửa văn phòng lập tức rung lắc dữ dội, vài tờ thậm chí bị gió xé tung, cuốn bay giữa không trung như những cánh chim vội vã tìm đường về tổ.

Tôi khom người, với tay nhặt lấy một tờ rơi trôi đến sát chân mình. Hóa ra là một tờ quảng cáo về một sự kiện đặc biệt sắp diễn ra. Trên mặt tờ giấy là một biểu tượng vô cùng quen thuộc, một hình bóng mà tôi đã quá thân thuộc từ lâu, từ cái thời mà mọi thứ còn đơn giản hơn, khi mà những buổi tập chung chỉ đơn thuần là những buổi tập chung.

Đó là nụ cười luôn tràn đầy năng lượng, một tinh thần mãi mãi tiến về phía trước, và quan trọng nhất\dots

Là hình ảnh đôi cánh rộng lớn; không phải đôi cánh của một thiên thần, mà là của một loài rồng huyền thoại.

Tôi vô thức siết chặt tờ giấy trong lòng bàn tay, các ngón tay bóp chặt đến mức các mép giấy nhàu nát dưới sức nóng của lòng bàn tay. Nhưng đến cuối cùng, tôi vẫn chẳng thể hiểu nổi vì sao trong lòng lại dâng lên một nỗi bất an khó tả. Có gì đó đang ở phía trước, một cái gì đó mà tôi không thể đoán trước được, dù rằng tôi đã quá quen với việc dự đoán những điều sắp đến.

Dù sao đi nữa\dots\ mùa hè cũng đang đến rất gần\dots

Tháng Sáu sắp gõ cửa.

Một mùa hè mang dấu ấn đặc biệt.

Một mùa hè mà tôi có linh cảm rằng, sẽ không bao giờ giống với bất kỳ mùa hè nào chúng ta đã trải qua trước đây. Một mùa hè nóng bỏng, đầy những biến động bất ngờ\dots\ và cũng thật kỳ lạ. Mùa hè của năm COVID này, chúng ta sẽ cùng nhau làm những gì nhỉ? Những chuyến đi chơi biển vắng bóng du khách, những buổi tiệc ngoài trời với khoảng cách an toàn, hay những giờ phút chỉ đơn giản là được ngồi bên nhau mà chẳng phải lo nghĩ điều gì?

Đó là câu hỏi mà cả tôi và tất cả mọi người đều đang mong chờ có được câu trả lời.

Trong suốt ba tháng tiếp theo đây\dots

Từng chút một, mọi lời giải đáp sẽ dần lộ ra.

\end{document}